% BEMv18 appendix snippet
% Multi-knot certified-length exponent test.
% Alpha-blind: no observed fine-structure value is used.

\subsection{Multi-knot test for the certified-length exponent}

\paragraph{Status.}
This subsection describes a falsifier for the BEMv17 symbol claim that the
first surviving R--T residual is \(q_{-2}\).  If the residual has scale weight
\(L^{-2}\), then the finite correction should be normalized by
\[
\mathcal N_{\rm RT}=M_{\max}L_{\rm cert}^{2}.
\]

For each knot \(K_i\), define the raw pair correction
\[
\Delta F_i=\Delta F_{\rm pair}(K_i/0_1),
\]
the retained mode count \(M_i\), and the certified length \(L_i=L_{\rm cert}(K_i)\).
For a general monomial normalizer
\[
\mathcal N_{a,b}=M_i^a L_i^b,
\]
define
\[
R_{a,b}(K_i)=
\frac{|\Delta F_i|}
{M_i^aL_i^b}.
\]
The BEMv17 target is
\[
a=1,\qquad b=2.
\]

\paragraph{Diagnostic.}
The exponent \(b\) is tested by scanning \(R_{a,b}\) across multiple knots and
checking whether systematic correlation with \(L_i\) is reduced while the
finite correction remains subleading:
\[
\frac{R_{a,b}(K_i)}
{N_{\rm soft}V_{\rm soft}L_i}
\ll 1.
\]
This is not a proof, because knot-shape factors may also enter.  It is a
necessary falsifiability test: if \(b=2\) fails across several certified
lengths, the \(q_{-2}\) interpretation is weakened.

\paragraph{BEMv18 result.}
The selected exponent candidate is
\[
a=1.0,\qquad b=2.0,
\]
with formula
\[
M^1 L^2.
\]
The canonical \(q_{-2}\) candidate \(M_{\max}L_{\rm cert}^{2}\) has status
\[
\boxed{\text{PASS_CANONICAL_B2_SELECTED}}.
\]
